\section{Modifications to chapter 7}

The following modifications shall be applied to chapter 7 (\emph{Expressions}
[\texttt{expr}]).

\begin{modification}[Removal of 7.5.3/3]
\stdremove{If the expression this appears within the predicate of a contract assertion (6.11.1) (including as the result of an implicit transformation (7.5.5.1) and including in the bodies of nested
\emph{lambda-expressions}) and the
current class encloses the contract assertion, \texttt{const} is combined with
the \emph{cv-qualifier-seq} used to generate the
resulting type (see below).}
\end{modification}

\begin{modification}[Modification to 7.5.5.1/2]
If an id-expression \emph{E} denotes a non-static non-type member of some class
\emph{C} at a point where the current class (7.5.3) is \emph{X} and

\begin{itemize}[label=\textemdash]

\item
\emph{E} is potentially evaluated or \emph{C} is \emph{X} or a base class of
\emph{X}, and

\item
\emph{E} is not the \emph{id-expression} of a class member access expression (7.6.1.5), and

\item
\emph{E} is not the \emph{id-expression} of a \emph{reflect-expression} (7.6.2.10), and

\item
if \emph{E} is a \emph{qualified-id}, \emph{E} is not the un-parenthesized operand of
the unary \& operator (7.6.2.2),

\end{itemize}

the \emph{id-expression} is transformed into a class member access expression
using (\texttt{*this}) as the object expression.

\stdremove{If this transformation occurs in the predicate of a precondition assertion of a constructor of X or a postcondition assertion of a destructor of X, the expression is ill-formed.}
\newline

NOTE 2 If C is not X or a base class of X, the class member access expression is ill-formed. Also, if the id-expression occurs within a static or explicit object member function, the class member access is ill-formed.
\newline

This transformation does not apply in the template definition context (13.8.3.2).
\newline

\textbf{Note to editors}: \stdremove{EXAMPLE 1} shall be \textbf{completely removed}.
\end{modification}


\begin{modification}[Removal of 7.5.5.2/6]
\stdremove{Otherwise, if the \emph{unqualified-id} names a result binding
(9.4.2) attached to a function with return type \texttt{U},}

\begin{itemize}[label=\textemdash]

\item
\stdremove{if \texttt{U} is “reference to \texttt{T}”, then the type of the
expression is \texttt{const T};}

\item
\stdremove{otherwise, the type of the expression is \texttt{const U}.}

\end{itemize}

\end{modification}

\begin{modification}[Removal of 7.5.5.2/7]
\stdremove{Otherwise, if the \emph{unqualified-id} appears in the predicate of a contract
assertion \emph{C} (6.11) and the entity is}

\begin{itemize}[label=\textemdash]
\item
\stdremove{a variable declared outside of \emph{C} of object type \texttt{T},}

\item
\stdremove{a variable or template parameter declared outside of \emph{C} of type “reference
to \texttt{T}”, or}

\item
\stdremove{a structured binding of type \texttt{T} whose corresponding variable is declared
outside of \emph{C},
then the type of the expression is \texttt{const T}.}
\end{itemize}
\end{modification}

\begin{modification}[Removal of 7.5.5.2/8]
\textbf{Note to editors}: Example 2 shall be
\textbf{completely removed}.
\end{modification}

\begin{modification}[Modification of 7.5.5.2/12]
The expression is an xvalue if it is move-eligible (see below); an lvalue if the entity is a function, variable,
structured binding (9.7), \stdremove{result binding (9.4.2),} data member, 
or template parameter object; and a prvalue
otherwise (7.2.1); it is a bit-field if the identifier designates a bit-field.
\end{modification}

\begin{modification}[Removal of 7.5.5.2/13]
\stdremove{If an \emph{id-expression E} appears in the predicate of a function contract
assertion attached to a function \emph{f} and
denotes a function parameter of \emph{f} and the implementation introduces any temporary objects to hold the
value of that parameter as specified in 6.8.7,}

\begin{itemize}[label=\textemdash]
\item
\stdremove{if the contract assertion is a precondition assertion and the evaluation of the precondition assertion is
sequenced before the initialization of the parameter object, \emph{E} refers to the most recently initialized
such temporary object, and}

\item
\stdremove{if the contract assertion is a postcondition assertion, it is unspecified
whether \emph{E} refers to one of the
temporary objects or the parameter object; the choice is consistent within a single evaluation of a
postcondition assertion.}
\end{itemize}

\end{modification}

\begin{modification}[Removal of 7.5.5.2/14]
\stdremove{If an id-expression \emph{E} names a result binding in a postcondition assertion and the implementation introduces
any temporary objects to hold the result object as specified in 6.8.7, and the postcondition assertion is
sequenced before the initialization of the result object (7.6.1.3), \emph{E} refers to the most recently initialized such
temporary object.}
\end{modification}

\begin{modification}[Removal of 7.5.5.3/7]
\stdremove{If \emph{Q} appears in the predicate of a contract assertion \emph{C}
(6.11) and the entity is}

\begin{itemize}[label=\textemdash]
\item
\stdremove{a variable declared outside of \emph{C} of object type \texttt{T},}

\item
\stdremove{a variable declared outside of \emph{C} of type “reference to \texttt{T}”, or}

\item
\stdremove{a structured binding of type \texttt{T} whose corresponding variable is declared
outside of \emph{C},}

\end{itemize}

\stdremove{then the type of the expression is \texttt{const T}.}
\end{modification}

\begin{modification}[Modification of 7.5.5.3/8]
\stdremove{Otherwise, the} \stdadd{The} type of the expression is the type of the result.
\end{modification}

\begin{modification}[Modification to 7.5.6.1]
lambda-expression :\\
\tabto{1em}lambda-introducer attribute-specifier-seq\textsubscript{opt} lambda-declarator
compound-statement\\
\tabto{1em}lambda-introducer \textless template-parameter-list \textgreater
requires-clause\textsubscript{opt} attribute-specifier-seq\textsubscript{opt}\\
\tabto{1em}lambda-declarator compound-statement\\

lambda-introducer :\\
\tabto{1em}[ lambda-capture\textsubscript{opt} ]\\

lambda-declarator :
\tabto{1em}lambda-specifier-seq noexcept-specifier\textsubscript{opt}
attribute-specifier-seq\textsubscript{opt}
trailing-return-type\textsubscript{opt}\\
\tabto{2em}\stdremove{function-contract-specifier-seq\textsubscript{opt}}\\
\tabto{1em}noexcept-specifier attribute-specifier-seq\textsubscript{opt}
trailing-return-type\textsubscript{opt}
\stdremove{function-contract-specifier-seq\textsubscript{opt}}\\
\tabto{1em}trailing-return-type\textsubscript{opt}
\stdremove{function-contract-specifier-seq\textsubscript{opt}}\\
\tabto{1em}( parameter-declaration-clause ) lambda-specifier-seq\textsubscript{opt}
noexcept-specifier\textsubscript{opt}
attribute-specifier-seq\textsubscript{opt}\\
\tabto{2em}trailing-return-typeopt requires-clause\textsubscript{opt}
\stdremove{function-contract-specifier-seq\textsubscript{opt}}\\

lambda-specifier :
\tabto{1em}\texttt{consteval}\\
\tabto{1em}\texttt{constexpr}\\
\tabto{1em}\texttt{mutable}\\
\tabto{1em}\texttt{static}\\

lambda-specifier-seq :\\
\tabto{1em}lambda-specifier lambda-specifier-seq\textsubscript{opt}
\end{modification}

\begin{modification}[Modification to 7.5.6.2/7]
The function call operator or operator template is a static member function or static member function
template (11.4.9.2) if the \emph{lambda-expression}’s
\emph{parameter-declaration-clause} is followed by static. Otherwise,
it is a non-static member function or member function template (11.4.3) that is declared const (11.4.3) if
and only if the \emph{lambda-expression}’s \emph{parameter-declaration-clause}
is not followed by \texttt{mutable} and the lambda-
declarator does not contain an explicit object parameter. It is neither virtual nor declared volatile.
Any \emph{noexcept-specifier} \stdremove{or \emph{function-contract-specifier}
(9.4.1)} specified on a \emph{lambda-expression} applies to the
corresponding function call operator or operator template. An
\emph{attribute-specifier-seq} in a \emph{lambda-declarator}
appertains to the type of the corresponding function call operator or operator
template. An \emph{attribute-specifier-seq} in a \emph{lambda-expression}
preceding a \emph{lambda-declarator} appertains to the corresponding function
call operator or operator template. The function call operator or any given operator template specialization
is a \texttt{constexpr} function if either the corresponding
\emph{lambda-expression}’s \emph{parameter-declaration-clause} is followed
by \texttt{constexpr} or \texttt{consteval}, or it is \emph{constexpr-suitable} (9.2.6). It is an immediate function (9.2.6) if the
corresponding \emph{lambda-expression}’s \emph{parameter-declaration-clause} is
followed by \texttt{consteval}.
\end{modification}

\begin{modification}[Removal of 7.5.6.2/10]
\stdremove{If all potential references to a local entity implicitly captured by
a \emph{lambda-expression L} occur within the function
contract assertions (9.4.1) of the call operator or operator template of
\emph{L} or within \emph{assertion-statements} (8.9)
within the body of \emph{L}, the program is ill-formed.}
\newline

\stdremove{NOTE: 4 Adding a contract assertion to an existing C ++ program
cannot cause additional captures.}
\newline

\textbf{Note to Editors}: This modification also completely \textbf{removes the
corresponding code example}.
\end{modification}

\begin{modification}[Modification to 7.5.6.3/3]
A lambda-expression shall not have a \emph{capture-default} or
\emph{simple-capture} in its \emph{lambda-introducer} unless

\begin{itemize}[label=\textemdash]

\item
its innermost enclosing scope is a block scope (6.4.3),

\item
it appears within a default member initializer and its innermost enclosing scope is the corresponding
class scope (6.4.7), or

\item
\stdremove{it appears within a contract assertion and its innermost enclosing scope is the corresponding contract-
assertion scope (6.4.10).}

\end{itemize}

\end{modification}

\begin{modification}[Modification to 7.5.9/2.5]
Otherwise, if \emph{S} is a variable or a structured binding, \emph{S} shall either have static or thread storage duration
or shall inhabit a scope enclosing the expression. The expression \emph{E} is an lvalue referring to the object
or function \emph{X} associated with or referenced by \emph{S}, \stdadd{has the
same type as that of \emph{S},} and is a
bit-field if and only if \emph{X} is a bit-field. \stdremove{If \emph{E}
appears in the predicate of a contract assertion \emph{C} (6.11) and \emph{S}
is}

\begin{itemize}[label=\textemdash]

\item
\stdremove{a variable declared outside of \emph{C} of object type \texttt{T},}

\item
\stdremove{a variable declared outside of \emph{C} of type “reference to \texttt{T}”, or}

\item
\stdremove{a structured binding of type \texttt{T} whose corresponding variable is declared
outside of \emph{C},}
\end{itemize}

\stdremove{then the type of \emph{E} is \texttt{const T}, otherwise \emph{E} has
the same type as that of \emph{S}}.
\newline

NOTE: 2 The type of a \emph{splice-expression} designating a variable or structured binding of reference type will be
adjusted to a non-reference type (7.2.2).

\end{modification}

\begin{modification}[Modification of 7.6.1.3/6]
When a function is called, each parameter (9.3.4.6) is initialized (9.5, 11.4.5.3) with its corresponding
argument\stdremove{, and each precondition assertion of the function is
evaluated (9.4.1)}. If the function is an explicit
object member function and there is an implied object argument (12.2.2.2.2), the list of provided arguments is
preceded by the implied object argument for the purposes of this correspondence. If there is no corresponding
argument, the default argument for the parameter is used.\\
\ldots
\end{modification}

\begin{modification}[Modification of 7.6.1.3/7]
The \emph{postfix-expression} is sequenced before each expression in the
\emph{expression-list} and any default argument.
The initialization of a parameter or, if the implementation introduces any temporary objects to hold the
values of function parameters (6.8.7), the initialization of those temporaries, including every associated value
computation and side effect, is indeterminately sequenced with respect to that of any other parameter. 
\stdremove{These evaluations are sequenced before the evaluation of the precondition assertions of the function, which are
evaluated in sequence (9.4.1). For any temporaries introduced to hold the values of function parameters, the
initialization of the parameter objects from those temporaries is indeterminately sequenced with respect to
the evaluation of each precondition assertion.}
\end{modification}

\begin{modification}[Removal of 7.6.1.3/9]
\stdremove{When the called function exits normally (8.8.4, 7.6.2.4), all postcondition assertions of the function are
evaluated in sequence (9.4.1). If the implementation introduces any temporary objects to hold the result
value as specified in 6.8.7, the evaluation of each postcondition assertion is indeterminately sequenced with
respect to the initialization of any of those temporaries or the result object. These evaluations, in turn, are
sequenced before the destruction of any function parameters.}
\end{modification}
\begin{modification}[Modification to 7.6.2.4/2]
An \emph{await-expression} shall appear only as a potentially evaluated
expression within the \emph{compound-statement} of
a \emph{function-body} or \emph{lambda-expression}, in either case outside of a
handler (14.1). In a \emph{declaration-statement} or in
the \emph{simple-declaration} (if any) of an \emph{init-statement}, an
\emph{await-expression} shall appear only in an initializer of that
\emph{declaration-statement} or \emph{simple-declaration}. An
\emph{await-expression} shall not appear in a default argument (9.3.4.7).
An \emph{await-expression} shall not appear in the initializer of a block variable with static or thread storage
duration. \stdremove{An \emph{await-expression} shall not be a potentially-evaluated subexpression 
of the predicate of a contract
assertion (6.11).} A context within a function where an \emph{await-expression} can appear is called a suspension
context of the function.
\end{modification}

\begin{modification}[Modification to 7.7.4]
A variable \texttt{v} is \emph{constant-initializable} if

\begin{itemize}[label=\textemdash]
\item
the \emph{full-expression} of its initialization is a constant expression when
interpreted as a \emph{constant-expression}
\stdremove{with all contract assertions using the ignore evaluation semantic (6.11.2),}
\newline

NOTE 1 In the course of this determination,
\texttt{std::is\_constant\_evaluated()} (21.3.11) has the value \texttt{true}.
\newline

\stdremove{NOTE 2 Furthermore, if the initialization is manifestly
\emph{constant-evaluated}, its evaluation during translation can still evaluate contract assertions with other evaluation semantics, resulting in a diagnostic or ill-formed program if a contract violation occurs.}

\item
immediately after the initializing declaration of \texttt{v}, the object or
reference \texttt{x} declared by \texttt{v} is \emph{constexpr-representable}, and

\item
if \texttt{x} has static or thread storage duration, \texttt{x} is
\emph{constexpr-representable} at the nearest point whose
immediate scope is a namespace scope that follows the initializing declaration of v.

\end{itemize}

\end{modification}

\begin{modification}[Modification of 7.7.5/5]
\textbf{Note to editors}: \textbf{Remove} the following lines from example
\begin{verbatim}
constexpr int l(int c) pre(c >= 2) {
  return (c % 2 == 0) ? c / 0 : c;
}
const int i0 = l(0); // dynamic initialization; contract violation or undefined behavior
const int i1 = l(1); // static initialization; value of 1 or contract violation at compile time
const int i2 = l(2); // dynamic initialization; undefined behavior
const int i3 = l(3); // static initialization; value of 3
\end{verbatim}

\end{modification}

\begin{modification}[Modification to 7.7.7 Note 64]
64) Testing this condition can involve a notional evaluation of its
initializer\stdremove{, with evaluations of contract assertions using the
ignore evaluation semantic (6.11.2)}, as described above.
\end{modification}
